\documentclass[10pt,a4paper]{amsart}
\usepackage[margin=19mm]{geometry}
\usepackage{amsmath,amssymb,mathtools}
\usepackage[scr=rsfs,cal=euler]{mathalfa}
\usepackage{xfrac}
\usepackage[unicode,hidelinks,bookmarks=false]{hyperref}
\newcommand{\ie}{i.e.\ }
\newcommand{\cf}{cf.\ }
\newcommand{\resp}{respectively\ }
\newcommand{\Subsection}{\subsection}
\providecommand{\nobreakdash}{\nobreak\hbox{-}}
\setlength{\emergencystretch}{2em}
\multlinegap0pt
\usepackage{siunitx}
\usepackage{siunitx}
\usepackage{siunitx}
\usepackage{siunitx}
\usepackage{amssymb}
\usepackage{dsfont}
\usepackage{xcolor}
\usepackage{tikz}
\usepackage{cancel}
\usepackage{siunitx}
\newtheorem{prop}{}
\def\E{{\mathbb{E}}}
\def\P{{\mathbb{P}}}
\def\R{{\mathbb{R}}}
\title{Free boundary regularity and well-posedness of physical solutions to the supercooled Stefan problem}
\author{Sebastian Munoz}
\date{Source: arXiv:2506.18741v2 (CC BY 4.0)}
\begin{document}
\maketitle

\noindent\emph{Source and licence.} Free boundary regularity and well-posedness of physical solutions to the supercooled Stefan problem. Version arXiv:2506.18741v2 (CC BY 4.0), distributed by arXiv under the Creative Commons Attribution 4.0 International licence, http://creativecommons.org/licenses/by/4.0/, as recorded in the arXiv licence metadata for this exact version and shown on the abstract page https://arxiv.org/abs/2506.18741v2. Full licence notice: \texttt{licenses/arxiv-2506.18741-CC-BY-4.0.txt}.\par\smallskip

\noindent\textit{Editorial scope.} Counted unit: the article's prop eq (source0.tex lines 335--345). The article's complete original proof of it is delivered with it (source0.tex lines 346--380, one \texttt{proof} environment of 35 lines). No mathematics is written, repaired or re-derived: the statement and the proof are the article's own lines, in the article's own notation, and no retained line is substituted or re-flowed.  Retained uncounted as the source-local context the proof needs: lines 121--126; lines 146--147; lines 292--294; lines 301--303; lines 323--325.  Nothing was omitted from any retained span, so the package carries no omission record.  No retained line contains a citation, so the wrapper carries no bibliography and no bibliography entry.  No retained line contains a figure environment, an image-inclusion command or drawing code, so no image is delivered and the package declares no asset dependency.  Editorial formatting only, in addition: an AMS \texttt{amsart} preamble that restates the article's own declarations and macros used by the retained lines, namely the environments \texttt{prop} on separate counters, together with the standard packages those lines need.  The retained environments print in this document as Proposition 1 in order of appearance, whereas the article prints the same items with the numbering of its own full document.  The complete native source of the article as distributed in the arXiv e-print for this exact version is bundled unchanged as \texttt{sources/180358/source0.tex}.
\begin{equation}\begin{cases} \label{supercooled eqs}
    u_t - \frac{1}{2} u_{xx}= 0 & t>0,\; x>\Lambda_t,\\
\dot{\Lambda}_t = \frac{\alpha}{2} u_x(\Lambda_t,t) & t > 0,\\
u(x,0) = u_0(x),\; u(\Lambda_t,t) = 0 & x\geq 0,\; t\geq 0,
\end{cases}
\end{equation}

\begin{equation} \label{u defi}    \P(B_t\in [a,b], t<\tau)=\int_{a}^bu(x,t)dx,\quad \quad a<b,\; \;t>0,
\end{equation}

\begin{multline}\label{qwodkr8}
    0=\E(\varphi(B_{\tau},\tau))= \E(\varphi(B_0,0))+\E\left( \int_{0}^\tau \left(\frac{1}{2}\varphi_{xx}(B_s,s)+\varphi_t(B_s,s)\right)ds\right)
    \\=\E\left( \int_{0}^{\infty} \left(\frac{1}{2}\varphi_{xx}(B_s,s)+\varphi_t(B_s,s)\right)\chi_{s<\tau}ds\right)= \int_{0}^{\infty}\E\left(  \left(\frac{1}{2}\varphi_{xx}(B_s,s)+\varphi_t(B_s,s)\right),s<\tau\right)ds\\=\int_{0}^{\infty}\int_{\R} \left( \frac{1}{2}\varphi_{xx}+\varphi_t\right)u(x,s)dxds,\end{multline}

\begin{equation} \label{u(t-) formula}
    \P(B_{t_0}\in [a,b], t_0\leq \tau) =\int_{a}^bu(x,t_0-)dx,  \quad \Lambda_{t_0-} \leq a <b.
\end{equation}

  \begin{equation} \label{cov nu231}
    \int_{\Lambda_{s(x)}=\Lambda_{s(x)-}}\nu(x)dx=\alpha^{-1}\int_{\Lambda_{s(x)}=\Lambda_{s(x)-}}dx=\alpha^{-1}\int_{\Lambda_{t}=\Lambda_{t-}}d\Lambda_t.
    \end{equation}

\begin{prop} \label{prop: u eq} Let $\Lambda$ be a probabilistic solution to \eqref{supercooled eqs}, and let $u$ be given by \eqref{u defi}. The function $u$ satisfies, in the distributional sense,
    \begin{equation}\label{u equation}\begin{cases} u_t-\frac12 u_{xx}=(\nu(x)\chi_{u>0})_t & x\in \R \times (0,\infty), \\
    u(x,0)=u_0(x) \chi_{x>\Lambda_0} &x\in \R.
    \end{cases}
\end{equation}
Namely, for every $\varphi\in C_c^{\infty}(\R \times [0,\infty))$,
\begin{equation}\label{u equation dist}
     \int_{0}^{\infty}\int_{\R} \left( \varphi_t(u-\nu \chi_{u>0})+\frac{1}{2}\varphi_{xx}u\right)dxdt+\int_{\R}(u_0(x)-\nu(x))\varphi(x,0)dx=0.
\end{equation}
In particular, $u$ is a global subsolution to the heat equation in $\R \times (0,\infty)$, and $u \in L^2_{\operatorname{loc}}((0,\infty);H^1_{\operatorname{loc}}(\R))$.
\end{prop}

\begin{proof} Let $\varphi\in C_c^{\infty}(\R \times [0,\infty))$. Using It\^o's formula as in \eqref{qwodkr8},
\begin{multline}\label{qwodkr8u392qwxk12xv}
    \E(\varphi(B_{\tau},\tau))= \E(\varphi(B_0,0))+\E\left( \int_{0}^\tau \left(\frac{1}{2}\varphi_{xx}(B_s,s)+\varphi_t(B_s,s)\right)ds\right)
    \\=\int_{\R}u_0(x)\varphi(x,0)dx+ \int_{0}^{\infty}\int_{\R} \left( \frac{1}{2}\varphi_{xx}+\varphi_t\right)u(x,t)dxdt.\end{multline}
    Since $B_{\tau}=\Lambda_{\tau}$ whenever $\Lambda_{\tau}=\Lambda_{\tau^-}$, we have
\begin{equation}  \label{expcqwsum02s}\E(\varphi(B_{\tau},\tau))=\E(\varphi(\Lambda_{\tau},\tau), \Lambda_{\tau}=\Lambda_{\tau^-})+\E(\varphi(B_{\tau},\tau),\Lambda_{\tau}>\Lambda_{\tau^-}).\end{equation}
Now, observing that $\alpha^{-1}\Lambda_t=\P(\tau \leq t)$ is the cumulative distribution function of $\tau$, it follows that
    \begin{equation} \label{qwodkpsxk12xv}
        \E(\varphi(\Lambda_{\tau},\tau), \Lambda_{\tau}=\Lambda_{\tau^-})=\int_{\Lambda_s=\Lambda_{s-}}\varphi(\Lambda_s,s)d(\alpha^{-1}\Lambda_s).
    \end{equation}
Therefore, changing variables as in \eqref{cov nu231}, we have
     \begin{multline} \label{0392kcmew0mc0m}
         \int_{\Lambda_s=\Lambda_{s-}}\varphi(\Lambda_s,s)d(\alpha^{-1}\Lambda_s)
         =\alpha^{-1}\int_{{\Lambda_{s(x)}=\Lambda_{s(x)-}}}\varphi( x,s(x))dx\\=\alpha^{-1}\int_{\Lambda_{s(x)}=\Lambda_{s(x)-}}\int_{0}^{s(x)}\varphi_tdtdx+\alpha^{-1}\int_{\Lambda_{s(x)}=\Lambda_{s(x)-}}\varphi(x,0)dx\\=\int_{\Lambda_{s(x)}=\Lambda_{s(x)-}}\int_{0}^{\infty}\varphi_t\nu\chi_{u>0}dtdx+\int_{\Lambda_{s(x)}=\Lambda_{s(x)-}}\varphi(x,0)\nu(x)dx,
     \end{multline}
where the integrals are understood to be over $x\in (0,\alpha)$. Thus, we infer from \eqref{qwodkpsxk12xv} and \eqref{0392kcmew0mc0m} that
\begin{equation} \label{epx1dwlpas}
    \E(\varphi(\Lambda_{\tau},\tau), \Lambda_{\tau}=\Lambda_{\tau^-})=\int_{\Lambda_{s(x)}=\Lambda_{s(x)-}}\int_{0}^{\infty}\varphi_t\nu\chi_{u>0}dtdx+\int_{\Lambda_{s(x)}=\Lambda_{s(x)-}}\varphi(x,0)\nu(x)dx.
\end{equation}
On the other hand, since the set of discontinuities of $\Lambda$ is countable, using \eqref{u(t-) formula} and recalling the convention $u(\cdot,0-):=u_0$, we have
\begin{multline}  \label{qw2dklcfpqws2}   \E(\varphi(B_{\tau},\tau),\Lambda_{\tau}>\Lambda_{\tau^-})
=\sum_{\Lambda_{t}>\Lambda_{t-}}\E(\varphi(B_{\tau},\tau),\tau= t)
=\sum_{\Lambda_{t}>\Lambda_{t-}}\E(\chi_{B_t \in [\Lambda_{t-},\Lambda_t]}\varphi(B_{t},t),\tau\geq t)\\=\sum_{\Lambda_{t}>\Lambda_{t-}}\int_{\Lambda_{t-}}^{\Lambda_{t}}\varphi(x,t)u(x,t-)dx
=\sum_{\Lambda_{t}>\Lambda_{t-}}\int_{\Lambda_{t-}}^{\Lambda_{t}}\varphi(x,t)\nu(x)dx
\\=\sum_{\Lambda_{t}>\Lambda_{t-}}\int_{\Lambda_{t-}}^{\Lambda_{t}}\int_{0}^{\infty}\varphi_t\nu\chi_{u>0}dtdx+\int_{\Lambda_{t-}}^{\Lambda_t}\varphi(x,0)\nu(x)dx\\
=\int_{\Lambda_{s(x)}>\Lambda_{s(x)-}}\int_0^{\infty}\varphi_t\nu \chi_{u>0}dtdx+\int_{\Lambda_{s(x)}>\Lambda_{s(x)-}}\varphi(x,0)\nu(x)dx.
\end{multline}
In view of \eqref{expcqwsum02s}, \eqref{epx1dwlpas}, and \eqref{qw2dklcfpqws2}, using the fact that $\varphi$ is compactly supported and $\nu$ is supported in $[0,\alpha]$, we have
\begin{equation}
\E(\varphi(B_{\tau},\tau))=\int_{0}^{\infty}\int_{\R}\varphi_t(x,t)\nu(x)\chi_{u>0}dxdt+\int_{\R}\varphi(x,0)\nu(x)dx. 
\end{equation}
The equality \eqref{u equation dist} thus follows from \eqref{qwodkr8u392qwxk12xv}. Finally, the fact that $u$ is subcaloric follows readily from the fact that the set $\{u(t)>0\}$ is decreasing in time. Namely,
\[u_t-\frac12u_{xx}=(\nu\chi_{u>0})_t\leq0,\]
and the $L^2_{\operatorname{loc}}((0,\infty);H^1_{\operatorname{loc}}(\R))$ bounds on $u$ then follow from this inequality by standard energy estimates.
\end{proof}

\end{document}
